\documentclass[10pt,a4paper]{amsart}
\usepackage[margin=19mm]{geometry}
\usepackage{amsmath,amssymb,mathtools}
\usepackage[scr=rsfs,cal=euler]{mathalfa}
\usepackage{xfrac}
\usepackage[unicode,hidelinks,bookmarks=false]{hyperref}
\newcommand{\ie}{i.e.\ }
\newcommand{\cf}{cf.\ }
\newcommand{\resp}{respectively\ }
\newcommand{\Subsection}{\subsection}
\providecommand{\nobreakdash}{\nobreak\hbox{-}}
\setlength{\emergencystretch}{2em}
\multlinegap0pt
\usepackage{bm}
\usepackage{bbm}
\usepackage{dsfont}
\usepackage{xspace}
\usepackage{cancel}
\usepackage{mathtools}
\usepackage{amsthm}
\makeatletter
\@ifundefined{theorem}{\newtheorem{theorem}{Theorem}}{}
\@ifundefined{lemma}{\newtheorem{lemma}[theorem]{Lemma}}{}
\makeatother
\title[Homogeneous hypersurfaces of the four-dimensional Thurston g]{Homogeneous hypersurfaces of the four-dimensional Thurston geometry ${\rm Sol\_0^4}$}
\author{Marie D'haene, Guoxin Wei, Zeke Yao, Xi Zhang}
\date{Source: arXiv:2508.10545v2 (CC BY 4.0)}
\begin{document}
\maketitle

\noindent\emph{Source and licence.} Homogeneous hypersurfaces of the four-dimensional Thurston geometry ${\rm Sol_0^4}$. Version arXiv:2508.10545v2 (CC BY 4.0), distributed under the Creative Commons Attribution 4.0 International licence, as recorded in the arXiv licence metadata for this exact version and shown on the abstract page https://arxiv.org/abs/2508.10545v2. Full rights notice: licenses/arxiv-2508.10545-CC-BY-4.0.txt.\par\smallskip

\noindent\textit{Editorial scope.} Counted unit: the Lemma whose source label is \texttt{lem:4.4} in the article (source0.tex lines 1166--1169, source label \texttt{lem:4.4}), delivered together with the article's own complete original proof of it (source0.tex lines 1171--1345, one \texttt{proof} environment of 175 lines). No mathematics is written, repaired or re-derived: statement and proof are the article's own text line for line. Required context retained uncounted: eqn:2.2 (lines 334--339); eqn:2.16 (lines 932--936); eqn:2.20 (lines 940--946); eqn:2.18 (lines 965--970); eqn:8.20 (lines 974--976); eqn:9.45 (lines 980--986). Every definition, display and label of those lines is retained unchanged. Omitted from this package and not redistributed: 1--333, 340--931, 937--939, 947--964, 971--973, 977--979, 987--1165, 1170--1170, 1346--1890. Nothing inside a retained excerpt is omitted or altered: every retained body is delivered exactly as the article's lines are written, so no omission receipt of any kind accompanies this package. Citations: the retained excerpts carry no citation command at all, so no bibliography entry of the article is transcribed and no reference is redistributed. Reference adaptation: none was needed and none was made; every reference inside the delivered unit resolves inside it, so no unresolved reference or citation is produced. Printed slips kept literal: no printed word, symbol, hypothesis or number of the counted statement or of its proof was altered. Layout and class: the article's own class is \texttt{amsart} with packages including mathrsfs, amsfonts, amsmath, amssymb, amscd, amsbsy, bm, amsthm, enumerate, babel, hyperref, geometry, xcolor, caption; the wrapper is the lane's pinned \texttt{amsart} editorial wrapper at 10pt on A4, which restates the theorem-like environments the retained text needs and adds the article's own macro definitions that the retained text needs (none), with extra wrapper packages (bm, bbm, dsfont, xspace, cancel, mathtools). The delivered numbering is the unit's own. Assets: no retained span contains a figure environment, a graphics-inclusion command, a PDF-page inclusion, a table or a TikZ picture; no image file is copied into this package, the package declares no asset dependency and no image file is redistributed with it. Licence: version arXiv:2508.10545v2 of this article is distributed under the Creative Commons Attribution 4.0 International licence, as recorded in the arXiv licence metadata for this exact version and shown on the abstract page https://arxiv.org/abs/2508.10545v2. A full rights notice is delivered at licenses/arxiv-2508.10545-CC-BY-4.0.txt. Characters: every retained excerpt is pure ASCII, so the delivered engine, XeLaTeX with its default Latin Modern text font, reports no missing character.
	\begin{equation}\label{eqn:2.2}
		E_1=e^t \partial_x, \quad
		E_2=e^t \partial_y,\quad
		E_3=e^{-2t} \partial_z, \quad
		E_4= \partial_t.
	\end{equation}

	\begin{equation}\label{eqn:2.16}
		W_1=\frac{bE_1-a E_2}{\sqrt{1-d^2}}, \quad
		W_2=\frac{adE_1+bd E_2}{\sqrt{1-d^2}}-\sqrt{1-d^2}E_4, \quad
		W_3=E_3,
	\end{equation}

	\begin{equation}\label{eqn:2.20}
		\begin{aligned}
			&J_1W_1=J_2W_1=dW_2+\sqrt{1-d^2}N, \quad J_1W_2=-dW_1+\sqrt{1-d^2}W_3,\\
			&J_2W_2=-dW_1-\sqrt{1-d^2}W_3, \quad J_1W_3=-J_2W_3=-\sqrt{1-d^2}W_2+dN,\\
			& PW_1=0, \quad PW_2=(1-d^2)W_2-d\sqrt{1-d^2}N, \quad PW_3=0.
		\end{aligned}
	\end{equation}

		\begin{equation}\label{eqn:2.18}
			W_1(a)=-\frac{b\sqrt{1-d^2}}{d}, \quad 
			W_1(b)=\frac{a\sqrt{1-d^2}}{d},\quad 
			W_2(a)=W_3(a)
			= W_2(b)=W_3(b)=0.
		\end{equation}

		\begin{equation}\label{eqn:8.20}
			AW_1=\tfrac{1}{d}W_1, \quad AW_2=dW_2, \quad AW_3=-2dW_3.
		\end{equation}

		\begin{equation}\label{eqn:9.45}
			\begin{aligned}
				\nabla_{W_i}W_j&=0 \ ( i=1, 2, \ j=1, 2, 3), \quad \nabla_{W_3}W_1=0,\\
				\nabla_{W_3}W_2&=-2\sqrt{1-d^2}W_3, \quad
				\nabla_{W_3}W_3=2\sqrt{1-d^2}W_2.
			\end{aligned}
		\end{equation}

	\begin{lemma}\label{lem:4.4}
		Assume that {\bf Case I-(i)} occurs, then up to isometries of ${\rm Sol_0^4}$,
		$M$ is an open part of $M_{1,r}$ for some $r>0$.
	\end{lemma}

	\begin{proof}
		On $\Omega_{1}$, it can be directly checked that the Gauss and Codazzi equations are satisfied by using the
		local orthonormal frame field $\{W_i\}_{i=1}^3$ from~\eqref{eqn:2.16}, together with \eqref{eqn:2.20}, \eqref{eqn:8.20}
		and \eqref{eqn:9.45}. 
		
		By using \eqref{eqn:2.2} and \eqref{eqn:2.16}, we get
		%
		\begin{equation}\label{eqn:2.23}
			W_1(t)=0, \quad W_2(t)=-\sqrt{1-d^2}, \quad W_3(t)=0,
		\end{equation}
		where $t$ is the fourth Cartesian coordinate function on $\mathbb{R}^4$.
		From \eqref{eqn:9.45} and $[W_i,W_j]=\nabla_{W_i}W_j-\nabla_{W_j}W_i$, we have
		%
		\begin{equation}\label{eqn:2.223}
			[W_1,W_2]=0, \quad [W_1,W_3]=0, \quad [W_2,W_3]=2\sqrt{1-d^2}W_3.
		\end{equation}
		%
		
		Now, we define a new frame field $\{X_i\}_{i=1}^3$ on $M$ as follows:
		%
		\begin{equation}\label{eqn:2.34}
			X_1=W_1, \quad X_2=W_2, \quad X_3=e^{2t}W_3.
		\end{equation}
		%
		By using \eqref{eqn:2.23}--\eqref{eqn:2.34},
		we can verify that $[X_i,X_j]=0$, $1\leq i<j\leq3$.
		Moreover, it follows from \eqref{eqn:2.18} and \eqref{eqn:2.34} that
		%
		\begin{equation}
			\label{eqn:2.31}    
			X_1(a)=-\frac{b\sqrt{1-d^2}}{d}, \quad
			X_1(b)=\frac{a\sqrt{1-d^2}}{d}, \quad
			X_2(a)=X_3(a)=X_2(b)=X_3(b)=0.
		\end{equation}
		Since the vector fields $\{X_i\}_{i=1}^3$ commute and are linearly independent, we can locally identify $M$ with an open subset $\Omega$ of $\mathbb{R}^3$
		and express the hypersurface $M$ by an immersion
		%
		\begin{equation*}
			\begin{aligned}
				\Phi:\Omega\subset\mathbb{R}^3 \longrightarrow {\rm Sol_0^4}:
				(u,v,w) \longmapsto (x(u,v,w),y(u,v,w),z(u,v,w),t(u,v,w)),
			\end{aligned}
		\end{equation*}
		%
		such that
		%
		\begin{equation}\label{eqn:2.36}
			\begin{aligned}
				d\Phi(\partial_u)
				&=(x_u,y_u,
				z_u,t_u)=X_1,\\
				d\Phi(\partial_v)
				&=(x_v,y_v,
				z_v,t_v)=X_2,\\
				d\Phi(\partial_w)
				&=(x_w,y_w,
				z_w,t_w)=X_3.
			\end{aligned}
		\end{equation}
		%
		
		From \eqref{eqn:2.31} and \eqref{eqn:2.36}, we know
		that $a$ and $b$ depend only on $u$ and
		%
		\begin{equation}\label{eqn:2.38}
			a_u=\frac{-b\sqrt{1-d^2}}{d}, \quad b_u=\frac{a\sqrt{1-d^2}}{d}.
		\end{equation}
		%
		Solving the equation \eqref{eqn:2.38}, we obtain
		%
		\begin{equation}\label{eqn:2.39}
			\begin{aligned}
				a(u)=c_1\cos\left(\tfrac{\sqrt{1-d^2}}{d}u\right)+c_2\sin\left(\tfrac{\sqrt{1-d^2}}{d}u\right), \quad
				b(u)=c_1\sin\left(\tfrac{\sqrt{1-d^2}}{d}u\right)-c_2\cos\left(\tfrac{\sqrt{1-d^2}}{d}u\right),
			\end{aligned}
		\end{equation}
		%
		where $c_1^2+c_2^2=1-d^2$ and $c_1, c_2$ are constants.
		
		According to \eqref{eqn:2.2}, \eqref{eqn:2.16} and \eqref{eqn:2.34}, we have
		%
		\begin{equation}\label{eqn:2.37}
			X_1= \left(\frac{be^t}{\sqrt{1-d^2}},-\frac{ae^t}{\sqrt{1-d^2}},0,0\right),\quad
			X_2= \left(\frac{ade^t}{\sqrt{1-d^2}},\frac{bde^t}{\sqrt{1-d^2}},0,-\sqrt{1-d^2}\right),\quad
			X_3= \left(0,0,1,0\right).
		\end{equation}
		%
		The equations \eqref{eqn:2.36} and \eqref{eqn:2.37} show that
		$x$ and $y$ depend only on $(u,v)$, $z$ depends only on $w$,
		and $t$ depends only on $v$.
		It also holds that $t_v=-\sqrt{1-d^2}$, i.e.\ $t(v)=-\sqrt{1-d^2}v+c_3$, where $c_3$ is a constant.
		Hence, $v(t)=-\tfrac{t-c_3}{\sqrt{1-d^2}}$ and
		%
		\begin{equation}\label{eqn:2.40}
			v_t=-\frac{1}{\sqrt{1-d^2}}.
		\end{equation}
		%
		In what follows we use $t$ instead of $v$ as a local coordinate.
		
		From \eqref{eqn:2.36} and \eqref{eqn:2.39}--\eqref{eqn:2.40}, we have
		%
		\begin{equation*}
			\begin{aligned}
				x_u
				&=\frac{e^t}{\sqrt{1-d^2}}
				\left(c_1\sin\left(\tfrac{\sqrt{1-d^2}}{d}u\right)-c_2\cos\left(\tfrac{\sqrt{1-d^2}}{d}u\right)\right),\\
				y_u
				&=\frac{-e^t}{\sqrt{1-d^2}}
				\left(c_1\cos\left(\tfrac{\sqrt{1-d^2}}{d}u\right)+c_2\sin\left(\tfrac{\sqrt{1-d^2}}{d}u\right)\right),\\
				x_t
				&=x_v v_t
				=\frac{-de^t}{1-d^2}
				\left(c_1\cos\left(\tfrac{\sqrt{1-d^2}}{d}u\right)+c_2\sin\left(\tfrac{\sqrt{1-d^2}}{d}u\right)\right),\\
				y_t
				&=y_v v_t
				=\frac{-de^t}{1-d^2}
				\left(c_1\sin\left(\tfrac{\sqrt{1-d^2}}{d}u\right)-c_2\cos\left(\tfrac{\sqrt{1-d^2}}{d}u\right)\right).
			\end{aligned}
		\end{equation*}
		%
		By solving these equations, we obtain
		%
		\begin{equation}\label{eqn:2.42}
			\begin{aligned}
				x(u,t) &= \frac{-de^t}{1-d^2}\left(c_1\cos\left(\tfrac{\sqrt{1-d^2}}{d}u\right)+c_2\sin\left(\tfrac{\sqrt{1-d^2}}{d}u\right)\right)+c_4,\\
				y(u,t) &= \frac{-de^t}{1-d^2}\left(c_1\sin\left(\tfrac{\sqrt{1-d^2}}{d}u\right)-c_2\cos\left(\tfrac{\sqrt{1-d^2}}{d}u\right)\right)+c_5,
			\end{aligned}
		\end{equation}
		%
		where $c_4, c_5$ are constants.
		
		From \eqref{eqn:2.36} and \eqref{eqn:2.37}, we have $z_w=1$, i.e.\ 
		$z=w+c_6$, where $c_6$ is a constant.
		Now, by reparametrizing to use $z$ as a local coordinate, we have
		%
		\begin{equation*}
			\begin{aligned}
				\Phi:\Omega\subset\mathbb{R}^3 \longrightarrow {\rm Sol_0^4}:
				(u,t,z)\longmapsto (x(u,t),y(u,t),z,t),
			\end{aligned}
		\end{equation*}
		%
		where $x(u,t), y(u,t)$ are given by \eqref{eqn:2.42}.
		After a left translation by $(-c_4,-c_5,0,0)$,
		which is an isometry of ${\rm Sol_0^4}$, we obtain
		%
		\begin{equation*}
			\begin{aligned}
				\Phi(u,t,z)
				=\Bigg(&\frac{-de^t}{1-d^2}\left(c_1\cos\left(\tfrac{\sqrt{1-d^2}}{d}u\right)
				+c_2\sin\left(\tfrac{\sqrt{1-d^2}}{d}u\right)\right), \\
				&\quad \frac{-de^t}{1-d^2}\left(c_1\sin\left(\tfrac{\sqrt{1-d^2}}{d}u\right)
				-c_2\cos\left(\tfrac{\sqrt{1-d^2}}{d}u\right)\right)
				,z,t\Bigg).
			\end{aligned}
		\end{equation*}
		%
		
		Since $0<d<1$, we may assume that $d=\tanh r$ for some constant $r>0$.
		Let $L_1=
		\cosh r \left(\begin{smallmatrix}
			-c_1 & c_2 \\
			-c_2 & -c_1 \\
		\end{smallmatrix}\right)
		\in {\rm SO(2)}$.
		Then $L_1(\Phi(u,t,z))=(e^t\cos ({\rm csch}\, (r) u)\sinh r,
		e^t\sin ({\rm csch}\, (r) u)\sinh r,z,t)$.
		By taking the reparametrization $x_1={\rm csch}\, (r) u$, $x_2=z$
		and $x_3=t-\ln({\rm sech}\, r)$, we get
		$$
		L_1(\Phi(x_1,x_2,x_3))=(e^{x_3}\cos x_1\tanh r,e^{x_3}\sin x_1\tanh r,
		x_2,x_3+\ln({\rm sech}\, r)).
		$$
		This shows that $M$ is an open part of $M_{1,r}$ for some $r>0$.
	\end{proof}

\end{document}
